\section{Introduction to the Topic and Historical Context}

\subsection{Introduction}

Throughout Latin America in 1975, the Cold War has become more than a confrontation between two superpowers: it has evolved into a struggle over the political future of the continent itself. Across the Southern Cone, revolutionary organizations, military governments, intelligence agencies, and foreign powers compete to shape the region according to fundamentally different political and ideological visions.

The success of the Cuban Revolution in 1959 transformed armed revolution into a viable political model for many left-wing movements across the continent. During the following years, organizations inspired by Marxism, socialism, nationalism, and anti-imperialism emerged in several countries, combining political mobilization with guerrilla warfare and clandestine organization. While these movements differ in ideology, leadership, and long-term objectives, many share a common commitment to resisting authoritarian rule and challenging the influence of the United States in Latin America.

Their expansion, however, is met by an equally determined response. Throughout the 1960s and early 1970s, military coups in Brazil, Bolivia, Uruguay, and Chile establish governments committes to counter-insurgency under the National Security Doctrine. Political violence intensifies across the region, exile communities continue to grow, and revolutionary organizations increasingly operate across national borders in search of security, resources, and opportunities for cooperation. At the same time, indications of closer contact between neighboring security and intelligence services suggest that state responses to these movements may also be becoming increasingly interconnected.

Against this backdrop, Latin America entrance a period of profound uncertainty. Revolutionary organizations continue to pursue their political objectives under growing pressure, military governments seek to consolidate their authority, and both the United States and the Soviet Union view the region as an increasingly important theatre of Cold War competition. The balance between revolution and counter-revolution remains unresolved, and the decisions taken by each of the actors involved have the potential to reshape the political future of the continent.

\subsection{Key Terms}

\begin{guidetable}[Y{0.52}Y{1.48}]{2}
  \textbf{Guerrilla Warfare} & Heavily informed by the Cuban Revolution and Che Guevara's writings, this operational doctrine rejects conventional confrontation in favor of tactical mobility, command decentralization, and deep underground activity. \newline Crucially, this framework functions as a political-military dual structure rather than a purely armed effort: while rural factions focus on territorial insurgencies, urban cells embed student and labor movements into covert networks. Operational methods—ranging from sabotage and intelligence to propaganda, kidnappings, and financial expropriations—ultimately serve a broader strategic imperative: political recruitment, mass mobilization, and ideological indoctrination. \\
  \midrule
  \textbf{Counterintelligence} & Under systemic military repression, counterintelligence mandate the implementation of rigorous operational security protocols. Its primary mechanism is structural compartmentalization, isolating information within self-contained cells, alongside communication security and the establishment of covert channels for sensitive intelligence data. \newline Rather than a secondary administrative task, counterintelligence is just as strategically vital as armed operations. Given the high-risk environment, a single intelligence breach does not merely jeopardize data; it routinely triggers the systematic dismantling of entire networks, mass arbitrary detentions, and the enforced disappearance of key operatives. \\
  \midrule
  \textbf{The National Security Doctrine} & Promoted under U.S. geopolitical influence, the doctrine pivots the military focus away from conventional territorial defense against foreign powers toward combating an "internal enemy”. \newline Under this framework, state security apparatuses cease to view trade unions, student movements, leftist political parties, and dissidents as legitimate political actors, recategorizing them instead as subversion threats to state survival. \newline Ultimately, the NSD serves as the primary ideological justification for state-sponsored terrorism. By viewing domestic politics through the lens of warfare, military regimes rationalize systematic human rights violations, such as censorship, arbitrary detentions, torture, and enforced disappearances, as necessary security measures to safeguard the nation. \\
  \midrule
  \textbf{State Terrorism} & Unlike conventional warfare, state terrorism relies on the directed use of state coercive apparatuses—military forces, police units, and intelligence agencies—to carry out extrajudicial executions, arbitrary arrests, systemic torture, and enforced disappearances against a nation's own civilian population. \newline Far from being isolated operational excesses, these actions constitute an organized state policy designed to neutralize revolutionary organizations and suppress civil resistance. Through various mechanisms, this strategy expands transnationally, facilitating cross-border intelligence sharing and targeted political assassinations of exiles. Ultimately, state terrorism aims not only at the physical elimination of political dissidents, but also at the psychological disruption of public morale and social solidarity. \\
  \midrule
  \textbf{Proxy Wars} & Following the 1959 Cuban Revolution, Latin America emerged as a central theater in this Cold War rivalry. For the United States, containing communism means channeling economic aid, military training, and intelligence sharing toward anti-communist regimes. Concurrently, the Soviet Union and Cuba may provide political, financial, and logistical backing to revolutionary organizations aligned with socialist ideals. \newline However, external assistance rarely arrives without strategic conditions. Superpowers frequently subordinate local objectives to their broader global agendas, creating a tension between ideological solidarity and strategic autonomy. Consequently, foreign aid must be evaluated critically to determine whether external support enhances revolutionary capabilities or merely fosters new structures of geopolitical dependency. \\
  \midrule
  \textbf{Southern Cone} & Mainly referring to Argentina, Chile, Uruguay, Paraguay and Brazil here. Most of these countries are governed by military or authoritarian regimes. \\
  \midrule
  \textbf{Internal Enemy} & Concept of the National Security Doctrine, that perceives the political opponents as threats from foreign military enemies. This term can refer to guerrillas, students, workers, activists… \\
\end{guidetable}

\subsection{Historical context}

\subsubsection{The Cold War in Latin America}

Following the 1959 Cuban Revolution, revolutionary movements gained momentum across Latin America, inspired by socialist, Marxist, and anti-imperialist ideologies. This new political landscape alarmed not only the United States military but also the governments of various Latin American states. As tensions escalated, many of these governments began to treat left-wing opposition groups not as legitimate political adversaries, but as illegal enemies.

The USA perceived the expansion of communist ideas within Latin America to disturb the security and position of their country. The Cuban Revolution inspired many left-wing movements showcasing the ability of overthrowing a government and establishing a social state.

The Soviet Union and Cuba supported the different revolutionary groups at the moment. But mainly these groups focused on national priorities more than working under Soviet or Cuban control. The conflicts escalated and the local revolutionary movements were supported internationally. By early 1970s there were armed organisations, guerrilla movements, labour movements, students, and political parties.

\subsubsection{Revolutionary Movements}

Before Operation Condor’s existence, revolutionary movements were already present in Latin America. Some groups were:

\begin{guidetable}[Y{0.42}Y{0.46}Y{2.12}]{3}
  \textbf{\emph{\textbf{Country}}} & \textbf{\emph{\textbf{Movement}}} & \textbf{\emph{\textbf{What they Do}}} \\
  \midrule
  Uruguay & Tupamaros (MLN-T) & A pioneering urban guerrilla movement specializing in high-impact operations: fund expropriations (bank robberies), political kidnappings, large-scale prison breaks, and armed propaganda through the temporary seizure of radio stations or towns. \\
  Argentina & ERP and Montoneros & \textbullet~\textbf{ERP:} Armed wing of the PRT, focused on military operations against state forces, attacks on military units, establishing a rural insurgent front in Tucumán, and the seizure of factories. \newline \textbullet~\textbf{Montoneros:} Peronist armed organization specializing in sabotage, attacks on military leaders, the kidnapping of corporate executives for financing purposes, and labor union mobilization. \\
  Bolivia & ELN & Born out of the guerrilla movement led by Che Guevara, it focused its activities on clandestine armed resistance, the reorganization of rural and urban guerrilla cells, and the mobilization of university and mining support networks. \\
  Chile & MIR & A Marxist–Leninist organization focused on land occupations for peasants and workers (estate seizures), the formation of popular militias, paramilitary training, and, following the 1973 coup, clandestine resistance and sabotage against the military regime. \\
\end{guidetable}

These challenged the authoritarian governments in different manners: armed struggle, political organising, propaganda, student organizations, exile support… They were connected by a common enemy but they weren’t always united.

\subsubsection{The Rise of Dictatorships}

By November 1975, the Southern Cone of Latin America had become a stronghold of authoritarian regimes. The rise of these dictatorships was not a series of isolated events, but rather a coordinated and planned regional counter-revolutionary strategy, heavily influenced by the Cold War and the National Security Doctrine.

This process began in 1960 and accelerated during the 1970s, as various countries suffered military coups that resulted in the systematic dismantling of democratic institutions.

\begin{itemize}
  \item \textbf{Paraguay (1954):} Alfredo Stroessner established a dictatorship that turned Paraguay into an early hub for cross-border intelligence exchange, political repression, and a safe haven for regional dictators.
  \item \textbf{Brazil (1964):} The overthrow of President João Goulart marked the first major military coup by an authoritarian regime during the Cold War in the region. The Brazilian military regime served as a model for neighboring countries, introducing state terrorism, institutionalized torture, and early international surveillance mechanisms.
  \item \textbf{Bolivia (1971):} Following years of political instability, General Hugo Banzer seized power through a bloody military coup, violently suppressing left-wing miners, peasant movements, and the remnants of the National Liberation Army (ELN).
  \item \textbf{Uruguay (1973):} Uruguay had historically been considered the "Switzerland of the Americas," but it fell under authoritarian rule when President Juan María Bordaberry, backed by the armed forces, dissolved parliament to crush the Tupamaros and suppress trade unions and student movements.
  \item \textbf{Chile (1973):} On September 11, General Augusto Pinochet ousted the democratically elected socialist President Salvador Allende. The regime immediately launched a brutal campaign of persecution against the Revolutionary Left Movement (MIR), the Socialist Party, and the Communist Party.
  \item \textbf{Argentina:} Although the government is still under Isabel Perón’s rule, state-sponsored cross-border repression and paramilitary violence, carried out notably by the Argentine Anticommunist Alliance (Triple A), is already fully operational, targeting the ERP, Montoneros, and foreign exiles.
\end{itemize}

\subsubsection{The Perception of the Dictatorships: The "Internal Enemy"}

By the early 1970s, Southern Cone military regimes no longer viewed political opponents through the lens of democratic rivalry. Driven by Cold War anti-communism and the National Security Doctrine, they framed dissent as an existential threat, defining any opposition as an "internal enemy" (\emph{enemigo interno}) rather than a legitimate political voice.

This shift has blurred the line between armed guerrillas and peaceful civilians. While insurgent groups remain primary targets, the "subversive" label has quickly expanded to swallow trade unionists, student organizers, journalists, academics, liberation theology priests, human rights advocates, lawyers, and even the relatives of suspects. Individuals are targeted less for specific criminal acts than for their ideological leanings, social circles, or potential to disrupt the state.

The entire repressive apparatus is transforming accordingly. Intelligence services have abandoned traditional criminal investigation to focus on dismantling perceived subversive networks. Surveillance, infiltration, arbitrary arrest, torture, forced disappearances, and extrajudicial killings have been rebranded as preventive tactics in an ongoing war. Legal protections and due process are dismissed as mere impediments to national security, clearing the way for state-sponsored violence against civilians.

For these regimes, the war is not just against guerrillas, but against the broader social fabric that sustains radical ideas.

\subsubsection{Exile and Cross-Border Networks}

Political exile has become one of the defining features of the Southern Cone during the 1970s. As military dictatorships consolidate their power, thousands of activists, union members, students, academics, journalists, and political leaders are forced to flee their countries for fear of imprisonment, torture, or death. Argentina, Chile, Uruguay, Bolivia, and Paraguay produce large communities of exiles who settle both in Latin America and abroad, particularly in Mexico, Venezuela, Cuba, and various European countries.

Exile, however, does not mark the end of revolutionary activity. On the contrary, it has given rise to a new dynamic, in which many resistance movements reorganize themselves across national borders, establishing clandestine communication channels, safe havens, financial networks, and international solidarity campaigns. Exile communities become essential for coordinating humanitarian aid, obtaining forged documents, recruiting new members, raising funds, disseminating information about human rights abuses, and maintaining contact with clandestine organizations operating in their countries of origin.

These transnational networks have soon become a major concern for military regimes. From their perspective, political refugees are not merely individuals fleeing persecution, but active participants in a broader revolutionary movement.
